\documentclass[a4paper,11pt]{article}
\usepackage[utf8]{inputenc}
\usepackage[french]{babel}
\usepackage[T1]{fontenc}
\usepackage{amsmath,amssymb,amsthm}
\usepackage{geometry}
\geometry{margin=2cm}
\usepackage{enumitem}
\usepackage{hyperref}
\usepackage{booktabs}
\usepackage{array}
\usepackage{xcolor}
\usepackage{listings}
\usepackage{titlesec}
\usepackage{fancyhdr}
\usepackage{lastpage}
\usepackage{graphicx}
\usepackage{etoolbox}
\AtBeginDocument{\shorthandoff{;:!?}}

\titleformat{\section}{\large\bfseries\color{blue!60!black}}{}{0pt}{}[\vspace{-0.5ex}\rule{\textwidth}{0.4pt}]
\titleformat{\subsection}{\bfseries\color{blue!40!black}}{}{0pt}{}
\titleformat{\subsubsection}{\itshape\color{blue!30!black}}{}{0pt}{}

\lstdefinestyle{monstyle}{
  frame=single,
  numbers=left,
  numbersep=5pt,
  breaklines=true,
  basicstyle=\small\ttfamily,
  backgroundcolor=\color{gray!5},
  rulecolor=\color{gray!40},
  columns=fullflexible,
  keepspaces=true,
  showstringspaces=false
}
\lstset{style=monstyle}

\pagestyle{fancy}
\fancyhf{}
\fancyhead[L]{\small STA211 -- Classification de variables}
\fancyhead[R]{\small Fiche r\'ecapitulative}
\fancyfoot[C]{\thepage/\pageref{LastPage}}
\renewcommand{\headrulewidth}{0.4pt}

\theoremstyle{definition}
\newtheorem*{definition}{D\'efinition}
\theoremstyle{remark}
\newtheorem*{remark}{Remarque}
\newtheorem*{important}{\`A retenir}

\title{\textbf{Fiche r\'ecapitulative -- STA211}\\[4pt]
\large Classification de variables}
\author{Vincent Audigier, Nd\`eye Niang}
\date{12 juin 2026}

\begin{document}
\maketitle
\thispagestyle{empty}

\begin{abstract}
Cette fiche pr\'esente les m\'ethodes de classification de variables.
Lorsque les variables sont nombreuses, il est utile de les regrouper en
classes homog\`enes pour analyser la multicolin\'earit\'e ou r\'eduire
la dimension. La fiche couvre les m\'ethodes hi\'erarchiques
(ascendante et descendante), les m\'ethodes de partitionnement direct
(CLV, ClustOfVar), les indices de similarit\'e selon la nature des
variables (quantitative, qualitative, binaire, mixte), et une
application sur les donn\'ees German Credit.
\end{abstract}

\vfill
\tableofcontents
\newpage

% ===================================================================
\section{Objectifs et contexte}
% ===================================================================

Quand les variables sont nombreuses, il est parfois tr\`es utile de
les regrouper en groupes homog\`enes. Ceci peut servir \`a~:
\begin{itemize}
  \item Analyser la multicolin\'earit\'e entre les variables.
  \item R\'eduire le nombre de variables en leur substituant un petit
    nombre de variables latentes synth\'etiques repr\'esentatives de
    chaque classe.
\end{itemize}

On distingue trois strat\'egies de partitionnement~:
\begin{enumerate}
  \item Les m\'ethodes hi\'erarchiques \textbf{ascendantes}.
  \item Les m\'ethodes hi\'erarchiques \textbf{descendantes}.
  \item Les m\'ethodes de \textbf{partitionnement direct}.
\end{enumerate}

% ===================================================================
\section{M\'ethodes hi\'erarchiques}
% ===================================================================

\subsection{M\'ethode ascendante (CAH)}

Les groupements se font par agglom\'eration progressive des variables
deux \`a deux. On r\'eunit les variables les plus proches au sens
d'une dissimilarit\'e, en utilisant un crit\`ere d'agr\'egation.

La complexit\'e est en $O(p^3)$ o\`u $p$ est le nombre de variables.

\subsubsection{Indices de similarit\'e/dissimilarit\'e}

Le choix de l'indice d\'epend de la nature des variables.

\paragraph*{Variables quantitatives~:}
le coefficient de corr\'elation lin\'eaire au carr\'e $r_{jj'}^2$
est l'indice naturel de similarit\'e. La dissimilarit\'e associ\'ee
$d = 1 - r^2$ est une distance euclidienne.

\paragraph*{Variables qualitatives~:}
les indices les plus classiques sont~:
\begin{itemize}
  \item $\chi^2$ de contingence
  \item Tschuprow ($T$)
  \item Cramer ($V$)
\end{itemize}

\paragraph*{Variables binaires~:}
plusieurs indices d\'efinis \`a partir du tableau de contingence
$2\times2$ ($N_{11}, N_{10}, N_{01}, N_{00}$)~:

\begin{tabular}{ll}
\toprule
Indice & Formule \\
\midrule
Sokal-Michener & $\displaystyle \frac{N_{11}+N_{00}}{N_{11}+N_{01}+N_{10}+N_{00}}$ \\
Jaccard & $\displaystyle \frac{N_{11}}{N_{11}+N_{01}+N_{10}}$ \\
Russel-Rao & $\displaystyle \frac{N_{11}}{N_{11}+N_{01}+N_{10}+N_{00}}$ \\
Ochiai & $\displaystyle \frac{N_{11}}{\sqrt{(N_{11}+N_{01})(N_{11}+N_{10})}}$ \\
Dice & $\displaystyle \frac{N_{11}}{2N_{11}+N_{01}+N_{10}}$ \\
Rogers-Tanimoto & $\displaystyle \frac{N_{11}+N_{00}}{N_{11}+2(N_{01}+N_{10})+N_{00}}$ \\
Kulzinsky & $\displaystyle \frac{N_{11}}{N_{01}+N_{10}}$ \\
\bottomrule
\end{tabular}

\begin{remark}
  Les mesures de dissimilarit\'e associ\'ees sont toutes euclidiennes
  (Fichet \& le Calve, 1984).
\end{remark}

\paragraph*{Variables mixtes~:}
deux m\'ethodes propos\'ees par Hummel, Edelmann \& Kopp-Schneider
(2017)~:
\begin{itemize}
  \item \textbf{CluMix-ama}~: similarit\'e bas\'ee sur la valeur
    absolue du $\gamma$ de Goodman-Kruskal pour les couples mixtes.
  \item \textbf{CluMix-dcor}~: similarit\'e bas\'ee sur une distance
    de Gower entre individus.
\end{itemize}
Ces m\'ethodes sont impl\'ement\'ees dans le package R \texttt{CluMix}
(archives CRAN).

\subsubsection{Crit\`eres d'agr\'egation}

Les crit\`eres sont les m\^emes que pour la CAH d'individus~:
\begin{itemize}
  \item Saut minimum (\emph{single linkage})
  \item Saut maximum (\emph{complete linkage})
  \item Saut moyen (\emph{average linkage})
  \item Ward (si la dissimilarit\'e est une distance euclidienne)
\end{itemize}

\subsection{M\'ethode descendante}

Proc\`ede par dichotomies successives pour construire un arbre
hi\'erarchique. L'interpr\'etation des classes est plus ais\'ee car
chaque classe est d\'ecrite par les r\`egles de divisions successives.

\paragraph*{M\'ethode VARCLUS (SAS)~:}
\begin{enumerate}
  \item R\'ealiser une ACP des variables.
  \item Si la seconde valeur propre $> 1$, retenir les deux premi\`eres
    composantes.
  \item Affecter chaque variable \`a la composante avec laquelle elle
    est le plus corr\'el\'ee.
  \item R\'ep\'eter sur chaque groupe tant que la seconde valeur
    propre $> 1$.
\end{enumerate}

\begin{important}
  La m\'ethode VARCLUS est diff\'erente de ClustOfVar. Il existe
  une impl\'ementation Python (\texttt{variable-clustering}) mais pas
  d'impl\'ementation R standard.
\end{important}

% ===================================================================
\section{Partitionnement direct}
% ===================================================================

\subsection{M\'ethode CLV (Vigneau \& Qannari, 2003)}

Destin\'ee aux variables quantitatives. Approche similaire aux
$K$-means, maximisant un crit\`ere de colin\'earit\'e intra-classe~:

\begin{equation}
\sum_{k=1}^K \sum_{j=1}^p \delta_{k,j} \; \text{cov}^2(x_j, u_k)
\quad\text{sous contrainte}\quad \|u_k\| = 1
\end{equation}

o\`u $\delta_{k,j}=1$ si la variable $j$ est dans la classe $k$, et
$u_k$ est la premi\`ere composante principale de la classe $k$.

Algorithme~:
\begin{enumerate}
  \item Initialiser les classes.
  \item \'Etape d'affectation~: calculer $\delta_{k,j}$ (affecter
    chaque variable \`a la classe dont le centro\"ide est le plus
    proche).
  \item \'Etape de mise \`a jour~: recalculer $u_k$ (1\`ere CP de
    chaque classe).
  \item R\'ep\'eter jusqu'\`a convergence.
\end{enumerate}

\subsection{ClustOfVar (Chavent et al., 2012)}

G\'en\'eralisation de CLV aux variables qualitatives ou mixtes~:
\begin{itemize}
  \item Utilise les composantes principales de l'ACM (qualitatives)
    ou de l'AFDM (mixtes) comme centro\"ide $u_k$.
  \item Similarit\'e mesur\'ee par le carr\'e du rapport de
    corr\'elation (qualitative) ou le carr\'e de la corr\'elation
    (quantitative).
\end{itemize}

\begin{important}
  Les algorithmes de partitionnement direct de variables restent
  sensibles \`a leur initialisation. Il convient de tester plusieurs
  initialisations et de retenir celle optimisant le crit\`ere.
\end{important}

% ===================================================================
\section{Application aux donn\'ees German Credit}
% ===================================================================

On utilise le jeu German Credit ($n=1000$, $p=21$, 7 quantitatives,
14 qualitatives) pour illustrer la classification de variables
mixtes.

\subsection{Chargement et importation}

\begin{lstlisting}[language=R]
library(CluMix); library(ClustOfVar)
library(cluster); library(FactoMineR)
library(DescTools)

don <- read.table(
  "https://archive.ics.uci.edu/ml/machine-learning-
   databases/statlog/german/german.data",
  sep = " ", stringsAsFactors = TRUE)
\end{lstlisting}

\subsection{CAH avec ClustOfVar}

\begin{lstlisting}[language=R]
var.factor <- which(sapply(don, class) == "factor")
var.numeric <- which(sapply(don, class) == "numeric" |
                     sapply(don, class) == "integer")
set.seed(1234)
res.cah.hclustvar <- hclustvar(X.quanti = don[,var.numeric],
                               X.quali = don[,var.factor])
k_hclustvar <- 4
res.cutree <- cutreevar(res.cah.hclustvar, k = k_hclustvar)
\end{lstlisting}

\subsection{CAH avec CluMix}

\begin{lstlisting}[language=R]
# CluMix-ama
res.dendro_ama <- dendro.variables(don)
res.dendro_ama <- as.hclust(res.dendro_ama)
k_ama <- 3; part.cah.CluMix_ama <- cutree(res.dendro_ama, k = k_ama)

# CluMix-dcor
res.dendro_dcor <- dendro.variables(don, method = "distcor")
res.dendro_dcor <- as.hclust(res.dendro_dcor)
k_dcor <- 4; part.cah.CluMix_dcor <- cutree(res.dendro_dcor, k = k_dcor)
\end{lstlisting}

\subsection{Partitionnement avec ClustOfVar}

\begin{lstlisting}[language=R]
res.kmeansvar <- kmeansvar(X.quanti = don[,var.numeric],
                           X.quali = don[,var.factor],
                           init = k_hclustvar, nstart = 100)

# Meilleure initialisation via CAH ou tirage aleatoire
myrandominit <- function(matdist, nb.centres){
  random_centre <- sample(seq.int(ncol(don)), size = nb.centres)
  random_part <- apply(matdist[,random_centre], 1, which.min)
  return(random_part)
}
matsim <- PairApply(don, mixedVarSim, symmetric = TRUE)
matdist <- as.matrix(sqrt(round(1 - matsim, 5)))
random_init <- replicate(100, myrandominit(matdist, k_hclustvar),
                         simplify = FALSE)

list.init <- list(cah.CluMix_dcor = part.cah.CluMix_dcor,
                  cah.hclustvar = part.cah.hclustvar)
list.init <- c(list.init, random_init)

res.kmeansvar.new <- sapply(list.init, FUN = kmeansvar,
  X.quanti = don[,var.numeric], X.quali = don[,var.factor],
  simplify = FALSE)
winner <- which.max(colSums(sapply(res.kmeansvar.new, "[[", "wss")))
res.kmeansvar <- res.kmeansvar.new[[winner]]
\end{lstlisting}

\subsection{Partitionnement avec PAM}

\begin{lstlisting}[language=R]
diss.CluMix_ama <- dist.variables(don)
res.pam_ama <- pam(diss.CluMix_ama, diss = TRUE, k = k_ama)

diss.CluMix_dcor <- dist.variables(don, method = "distcor")
res.pam_dcor <- pam(diss.CluMix_dcor, diss = TRUE, k = k_dcor)
\end{lstlisting}

\subsection{Comparaison des partitions}

Une ACM sur le tableau des partitions met en \'evidence quatre
groupes de variables~:

\begin{enumerate}
  \item \textbf{Profil personnel}~: Length.of.current.employment,
    Sex.Marital.Status, Duration.in.Current.address, Age.years,
    Housing, No.of.dependents.
  \item \textbf{Type de cr\'edit et profil financier}~: Duration,
    Purpose, Credit.Amount, Instalment.per.cent, Guarantors, Job,
    Telephone, Foreign.Worker.
  \item \textbf{Actifs du demandeur}~: Status, Savings account/bonds,
    Other.installment.plans, Creditability.
  \item \textbf{Historique}~: History, No.of.Credits.at.this.Bank.
\end{enumerate}

\begin{remark}
  La variable Creditability (r\'eponse) est li\'ee au groupe des
  actifs, ce qui sugg\`ere une relation privil\'egi\'ee avec ces
  variables.
\end{remark}

% ===================================================================
\section{Conclusion}
% ===================================================================

La classification de variables vise \`a regrouper les variables en
classes homog\`enes~: les variables d'un m\^eme groupe sont
fortement li\'ees, celles de groupes diff\'erents sont faiblement
li\'ees.

Applications importantes~:
\begin{itemize}
  \item \textbf{Fuzzy forests}~: variante des for\^ets al\'eatoires
    qui classifie d'abord les variables corr\'el\'ees, puis effectue
    une s\'election backward au sein de chaque groupe (package R
    \texttt{fuzzyforest}).
  \item \textbf{Biclustering} (co-clustering)~: classification
    simultan\'ee des individus et des variables pour identifier des
    groupes d'individus caract\'eris\'es par des sous-ensembles de
    variables.
\end{itemize}

% ===================================================================
\section*{Questions cl\'es (QCM)}
% ===================================================================

\addcontentsline{toc}{section}{Questions cl\'es (QCM)}

\begin{enumerate}

  \item Quel est l'objectif principal de la classification de
    variables~?
    \begin{enumerate}
      \item Classifier les individus en groupes homog\`enes
      \item Regrouper les variables en classes homog\`enes
      \item R\'eduire le nombre d'individus
      \item Augmenter la dimension des donn\'ees
    \end{enumerate}
    \textbf{R\'eponse~: b (regrouper les variables).}

  \item Quelle est la complexit\'e de la CAH de variables~?
    \begin{enumerate}
      \item $O(n^2)$
      \item $O(p^3)$
      \item $O(n)$
      \item $O(p)$
    \end{enumerate}
    \textbf{R\'eponse~: b ($O(p^3)$ o\`u $p$ est le nombre de
    variables).}

  \item Quel indice de similarit\'e est naturel pour des variables
    quantitatives~?
    \begin{enumerate}
      \item Le $\chi^2$
      \item Le V de Cramer
      \item La corr\'elation au carr\'e $r^2$
      \item Le Jaccard
    \end{enumerate}
    \textbf{R\'eponse~: c (coefficient de corr\'elation lin\'eaire
    au carr\'e).}

  \item Dans la m\'ethode VARCLUS, quel crit\`ere d\'eclenche une
    division~?
    \begin{enumerate}
      \item La premi\`ere valeur propre $> 1$
      \item La seconde valeur propre $> 1$
      \item Le nombre de variables $> 10$
      \item Le crit\`ere de Ward
    \end{enumerate}
    \textbf{R\'eponse~: b (seconde valeur propre de l'ACP $> 1$).}

  \item Quel crit\`ere est maximis\'e par la m\'ethode CLV~?
    \begin{enumerate}
      \item La variance inter-classe
      \item La colin\'earit\'e intra-classe $\sum \text{cov}^2(x_j, u_k)$
      \item Le $\chi^2$ de contingence
      \item La distance entre centro\"ides
    \end{enumerate}
    \textbf{R\'eponse~: b (somme des carr\'es des covariances entre
    les variables et le centro\"ide $u_k$).}

  \item Que repr\'esente $u_k$ dans la m\'ethode CLV~?
    \begin{enumerate}
      \item La moyenne des variables de la classe $k$
      \item La premi\`ere composante principale de la classe $k$
      \item Le m\'edo\"ide de la classe $k$
      \item La variance de la classe $k$
    \end{enumerate}
    \textbf{R\'eponse~: b (premi\`ere composante principale).}

  \item Comment ClustOfVar \'etend CLV aux donn\'ees mixtes~?
    \begin{enumerate}
      \item En utilisant la distance de Gower
      \item En utilisant l'AFDM pour d\'efinir les centro\"ides
      \item En supprimant les variables qualitatives
      \item En transformant les variables en binaires
    \end{enumerate}
    \textbf{R\'eponse~: b (AFDM pour les centro\"ides).}

  \item Quel package R impl\'emente ClustOfVar~?
    \begin{enumerate}
      \item cluster
      \item ClustOfVar
      \item CluMix
      \item FactoMineR
    \end{enumerate}
    \textbf{R\'eponse~: b (\texttt{ClustOfVar}).}

  \item Que sont les fuzzy forests~?
    \begin{enumerate}
      \item Des for\^ets al\'eatoires avec classification pr\'ealable
        des variables
      \item Des for\^ets al\'eatoires sans \'elagage
      \item Des for\^ets al\'eatoires sur donn\'ees floues
      \item Une m\'ethode de classification de variables
    \end{enumerate}
    \textbf{R\'eponse~: a (random forest avec classification
    pr\'ealable des variables corr\'el\'ees).}

  \item Qu'est-ce que le biclustering~?
    \begin{enumerate}
      \item Une classification d'individus seulement
      \item Une classification de variables seulement
      \item Une classification simultan\'ee d'individus et de
        variables
      \item Une m\'ethode de r\'egression
    \end{enumerate}
    \textbf{R\'eponse~: c (classification simultan\'ee individus
    $\times$ variables).}

\end{enumerate}

% ===================================================================
\begin{thebibliography}{9}
% ===================================================================

\bibitem{AudigierNiang2026}
  V.~Audigier, N.~Niang.
  \newblock \emph{STA211~: Classification de variables}.
  \newblock Cours CNAM, 23 f\'evrier 2026.

\bibitem{Chavent2012}
  M.~Chavent, V.~Kuentz-Simonet, B.~Liquet, J.~Saracco.
  \newblock \emph{ClustOfVar: An R Package for the Clustering of
    Variables}.
  \newblock J. Stat. Soft., 50(13):1--16, 2012.

\bibitem{VigneauQannari2003}
  V.~Vigneau, E.~M.~Qannari.
  \newblock \emph{Clustering of variables around latent components}.
  \newblock Comm. Stat. Simul. Comput., 32(4):1131--1150, 2003.

\bibitem{Hummel2017}
  M.~Hummel, D.~Edelmann, A.~Kopp-Schneider.
  \newblock \emph{Clustering of variables in the presence of mixed
    data}.
  \newblock Technical Report, 2017.

\bibitem{Saporta2006}
  G.~Saporta.
  \newblock \emph{Probabilit\'es, analyse des donn\'ees et
    statistique}.
  \newblock Technip, 2006.

\bibitem{FichetLeCalve1984}
  B.~Fichet, G.~le Calve.
  \newblock \emph{Structure g\'eom\'etrique des principaux indices
    de dissimilarit\'e sur signes de pr\'esence-absence}.
  \newblock Stat. Anal. Donn\'ees, 9(3):11--44, 1984.

\end{thebibliography}

\end{document}
